\begin{enumerate}

\item Se suma la resistencia térmica por área.

\[\fitmath{\mathrm{Res}=\frac{1}{\mathrm{hg}}+\frac{\mathrm{th}}{\mathrm{cond}}+\frac{1}{\mathrm{hc}}}\]

\[\fitmath{\mathrm{Res}=\frac{1}{(1540)}+\frac{(0.0012)}{(23)}+\frac{1}{(2320)}\approx 0.001132559\,\mathrm{m^{2}\,K/W}}\]

\item Se calcula el flujo estacionario.

\[\fitmath{\mathrm{flux}=\frac{\mathrm{Tg}-\mathrm{Tc}}{\frac{1}{\mathrm{hg}}+\frac{\mathrm{th}}{\mathrm{cond}}+\frac{1}{\mathrm{hc}}}}\]

\[\fitmath{\mathrm{flux}=\frac{(1580)-(720)}{\frac{1}{(1540)}+\frac{(0.0012)}{(23)}+\frac{1}{(2320)}}\approx 759342.3\,\mathrm{W/m^{2}}}\]

\item La cara caliente sigue del salto convectivo exterior.

\[\fitmath{\mathrm{Tw}=\mathrm{Tg}-\frac{\mathrm{Tg}-\mathrm{Tc}}{\mathrm{hg}\cdot \left(\frac{1}{\mathrm{hg}}+\frac{\mathrm{th}}{\mathrm{cond}}+\frac{1}{\mathrm{hc}}\right)}}\]

\[\fitmath{\mathrm{Tw}=(1580)-\frac{(1580)-(720)}{(1540)\cdot \left(\frac{1}{(1540)}+\frac{(0.0012)}{(23)}+\frac{1}{(2320)}\right)}\approx 1086.921\,\mathrm{K}}\]

\item La energía absorbida determina el mínimo de caudal.

\[\fitmath{\mathrm{mc}=\frac{\mathrm{area}\cdot \left(\mathrm{Tg}-\mathrm{Tc}\right)}{\left(\frac{1}{\mathrm{hg}}+\frac{\mathrm{th}}{\mathrm{cond}}+\frac{1}{\mathrm{hc}}\right)\cdot \mathrm{cp}\cdot \mathrm{dt}}}\]

\[\fitmath{\mathrm{mc}=\frac{(0.028)\cdot \left((1580)-(720)\right)}{\left(\frac{1}{(1540)}+\frac{(0.0012)}{(23)}+\frac{1}{(2320)}\right)\cdot (1080)\cdot (165)}\approx 0.119313\,\mathrm{kg/s}}\]

\end{enumerate}
